% Part: first-order-logic
% Chapter: completeness
% Section: identity

\documentclass[../../../include/open-logic-section]{subfiles}

\begin{document}

\olfileid{fol}{com}{ide}
\olsection{తాదాత్మ్యం}

\begin{explain}
ముందరి విభాగంలో ఇచ్చిన పద నమూనా నిర్మాణం, $\eq$ను కలిగి లేని
సమితులు~$\Gamma$కు ప్రథమక్రమ తర్కపు సంపూర్ణతను స్థాపించడానికి
సరిపోతుంది. $\eq$ను కలిగి లేని ప్రతి $!A \in \Gamma^*$ను (అందువల్ల
ప్రతి $!A \in \Gamma$ను కూడా) పద నమూనా సంతృప్తిపరుస్తుంది. కానీ
$\eq$ ఉంటే అది పనిచేయదు. ఎందుకంటే అప్పుడు~$\Gamma^*$లో
!!a{sentence}~$\eq[t][t']$ ఉండవచ్చు; పద నమూనాలో ఏ పదం విలువైనా ఆ
పదమే. అందువల్ల $t$, $t'$ వేర్వేరు పదాలైతే, పద నమూనాలో వాటి విలువలు
వరుసగా $t$, $t'$ వేర్వేరుగా ఉంటాయి; కనుక $\eq[t][t']$ అసత్యం.
అయితే ``వర్గీకరణ'' అనే నిర్మాణం ద్వారా దీన్ని సరిచేయవచ్చు.
\end{explain}

\begin{defn}
  $\Gamma^*$ అనేది~$\Lang L$లోని వాక్యాల అవైరుధ్య, !!{complete}
  సమితి అనుకుందాం. $\Lang L$లోని సంవృత పదాల సమితిపై $\approx$
  సంబంధాన్ని ఇలా నిర్వచిస్తాం:
  \[
  t \approx t' \text{\quad iff \quad} \eq[t][t'] \in \Gamma^*
  \]
\end{defn}

\begin{prop}
\ollabel{prop:approx-equiv}
$\approx$ సంబంధానికి కింది ధర్మాలు ఉన్నాయి:
\begin{enumerate}
\item $\approx$ స్వావర్తనమైనది.
\item $\approx$ సౌష్ఠవమైనది.
\item $\approx$ సంక్రామకమైనది.
\item $t \approx t'$, $f$ ఒక !!a{function}, ఇంకా $t_1$, \dots,
  $t_{i-1}$, $t_{i+1}$, \dots, $t_n$లు సంవృత పదాలైతే,
\[
\Atom{f}{t_1,\dots, t_{i-1}, t, t_{i+1}, \dots, t_n} \approx
\Atom{f}{t_1,\dots, t_{i-1}, t', t_{i+1}, \dots, t_n}.
\]
\item $t \approx t'$, $R$ ఒక !!a{predicate}, ఇంకా $t_1$, \dots,
  $t_{i-1}$, $t_{i+1}$, \dots, $t_n$లు సంవృత పదాలైతే,
\begin{multline*}
\Atom{R}{t_1,\dots, t_{i-1}, t, t_{i+1}, \dots, t_n} \in \Gamma^* \text{ iff } \\
\Atom{R}{t_1,\dots, t_{i-1}, t', t_{i+1}, \dots, t_n} \in \Gamma^*.
\end{multline*}
\end{enumerate}
\end{prop}

\begin{proof}
$\Gamma^*$ అవైరుధ్యమైనది, !!{complete} కనుక $\eq[t][t'] \in
\Gamma^*$ అయినప్పుడు, అప్పుడు మాత్రమే $\Gamma^* \Proves \eq[t][t']$.
అందువల్ల కింది వాటిని చూపితే చాలు:
\begin{enumerate}
\item అన్ని సంవృత పదాలు~$t$కు $\Gamma^* \Proves \eq[t][t]$.
\item $\Gamma^* \Proves \eq[t][t']$ అయితే $\Gamma^* \Proves
  \eq[t'][t]$.
\item $\Gamma^* \Proves \eq[t][t']$, $\Gamma^* \Proves
  \eq[t'][t'']$ అయితే $\Gamma^* \Proves \eq[t][t'']$.
\item $\Gamma^* \Proves \eq[t][t']$ అయితే,
\[
\Gamma^* \Proves
\eq[\Atom{f}{t_1,\dots,t_{i-1},t,t_{i+1},\dots,t_n}][\Atom{f}{t_1,\dots,t_{i-1},t',t_{i+1},\dots,t_n}]
\]
ప్రతి $n$-స్థాన !!{function}~$f$కు, సంవృత పదాలు $t_1$, \dots,
$t_{i-1}$, $t_{i+1}$, \dots,~$t_n$కు.
\sourcecorrection{OLTECOM-007}{ఆంగ్ల మూలంలో మొదటి ప్రమేయ పదంలో
$t_{i+1}$ తరువాత వరుసగా రెండు కామాలు ఉన్నాయి. పదాల జాబితాకు కావలసిన
ఒక్క కామాను మాత్రమే ఉంచాం.}
\item $\Gamma^* \Proves \eq[t][t']$, ఇంకా
$\Gamma^* \Proves
\Atom{R}{t_1,\dots,t_{i-1},t,t_{i+1},\dots,t_n}$ అయితే,
$\Gamma^* \Proves \Atom{R}{t_1,\dots,t_{i-1},t',t_{i+1},\dots,t_n}$;
ఇది ప్రతి $n$-స్థాన !!{predicate}~$R$కు, సంవృత పదాలు $t_1$, \dots,
$t_{i-1}$, $t_{i+1}$, \dots,~$t_n$కు వర్తిస్తుంది.
\end{enumerate}
\end{proof}

\begin{prob}
\olref[fol][com][ide]{prop:approx-equiv} నిరూపణను పూర్తి చేయండి.
\end{prob}

\begin{defn}
$\Gamma^*$ అనేది భాష~$\Lang L$లోని అవైరుధ్య, !!{complete} సమితి,
$t$ ఒక సంవృత పదం, $\approx$ ముందరి నిర్వచనంలో ఉన్నట్లే అనుకుందాం.
అప్పుడు:
\[
\equivrep{t}{\approx} = \Setabs{t'}{t'\in \Trm[L], t \approx t'}
\]
ఇంకా $\equivclass{\Trm[L]}{\approx} = \Setabs{\equivrep{t}{\approx}}{t \in \Trm[L]}$.
\end{defn}

\begin{defn}
\ollabel{defn:term-model-factor}
$\Struct M = \Struct M(\Gamma^*)$ అనేది \olref[mod]{defn:termmodel}
లోని~$\Gamma^*$ పద నమూనా అనుకుందాం. అప్పుడు
$\Struct{\equivclass{M}{\approx}}$ కింది !!{structure}:
\begin{enumerate}
\item $\Domain{\equivclass{M}{\approx}} = \equivclass{\Trm[L]}{\approx}$.
\item $\Assign{c}{\equivclass{M}{\approx}} = \equivrep{c}{\approx}$
\item $\Assign{f}{\equivclass{M}{\approx}}(\equivrep{t_1}{\approx}, \dots,
  \equivrep{t_n}{\approx}) = \equivrep{\Atom{f}{t_1,\dots, t_n}}{\approx}$
\item $\tuple{\equivrep{t_1}{\approx}, \dots, \equivrep{t_n}{\approx}} \in
  \Assign{R}{\equivclass{M}{\approx}}$ అయినప్పుడు, అప్పుడు మాత్రమే
  $\Sat{M}{\Atom{R}{t_1,\dots, t_n}}$; అంటే, అప్పుడు మాత్రమే
  $\Atom{R}{t_1,\dots, t_n} \in \Gamma^*$.
\end{enumerate}
\end{defn}

\begin{explain}
$\equivclass{\Trm[L]}{\approx}$లోని మూలకాల కోసం
$\Assign{f}{\equivclass{M}{\approx}}$, $\Assign{R}{\equivclass{M}{\approx}}$
లను $\equivrep{t}{\approx}$ అని సూచిస్తూ, అంటే
$t \in \equivrep{t}{\approx}$ అనే \emph{ప్రతినిధులు} ద్వారా
నిర్వచించామని గమనించండి. ఈ నిర్వచనాలు ఆ ప్రతినిధుల ఎంపికపై ఆధారపడవని
నిర్ధారించాలి; అంటే, అవే తుల్యతా వర్గాలను నిర్ణయించే మరొక ఎంపిక~$t'$
ఉంటే ($\equivrep{t}{\approx} = \equivrep{t'}{\approx}$), నిర్వచనాలు
అదే ఫలితాన్ని ఇవ్వాలి. ఉదాహరణకు, $R$ ఒక ఏకస్థాన !!{predicate} అయితే,
నిర్వచనపు చివరి ఖండం ప్రకారం $\equivrep{t}{\approx} \in
\Assign{R}{\equivclass{M}{\approx}}$ అయినప్పుడు, అప్పుడు మాత్రమే
$\Sat{M}{\Atom{R}{t}}$. $t \approx t'$ అయ్యే మరొక పదం~$t'$కు
$\Sat/{M}{\Atom{R}{t}}$ అయితే, నిర్వచనం $\equivrep{t'}{\approx}
\notin \Assign{R}{\equivclass{M}{\approx}}$ అని కోరుతుంది. $t \approx
t'$ అయితే $\equivrep{t}{\approx} = \equivrep{t'}{\approx}$; కానీ
$\equivrep{t}{\approx} \in \Assign{R}{\equivclass{M}{\approx}}$,
$\equivrep{t}{\approx} \notin \Assign{R}{\equivclass{M}{\approx}}$
రెండూ ఉండలేవు. అయితే ఇది జరగదని \olref{prop:approx-equiv} హామీ
ఇస్తుంది.
\end{explain}

\begin{prop}
$\Struct{\equivclass{M}{\approx}}$ సునిర్వచితమైనది. అంటే, $t_1$,
  \dots, $t_n$, $t_1'$, \dots, $t_n'$లు సంవృత పదాలు, ఇంకా
  $t_i \approx t_i'$ అయితే,
\begin{enumerate}
\item $\equivrep{\Atom{f}{t_1,\dots, t_n}}{\approx} =
    \equivrep{\Atom{f}{t_1',\dots, t_n'}}{\approx}$; అంటే,
  \[
  \Atom{f}{t_1,\dots, t_n} \approx \Atom{f}{t_1',\dots, t_n'}
  \]
  మరియు
\item $\Sat{M}{\Atom{R}{t_1,\dots, t_n}}$ అయినప్పుడు, అప్పుడు మాత్రమే
  $\Sat{M}{\Atom{R}{t_1',\dots, t_n'}}$; అంటే,
  \[
    \Atom{R}{t_1,\dots, t_n} \in \Gamma^* \text{ iff }
    \Atom{R}{t_1',\dots, t_n'} \in \Gamma^*.
  \]
\end{enumerate}
\end{prop}

\begin{proof}
~$n$పై ఆగమనం, \olref{prop:approx-equiv} ద్వారా అనుగమిస్తుంది.
\end{proof}

పద నమూనా సందర్భంలోలాగే, సత్య ఉపసిద్ధాంతాన్ని నిరూపించడానికి ముందు
కింది ఉపసిద్ధాంతం అవసరం.

\begin{lem}
\ollabel{lem:val-in-termmodel-factored} $\Struct M = \Struct M
 (\Gamma^*)$ అనుకుందాం. అప్పుడు $\Value{t}{\equivclass{M}{\approx}} = \equivrep{t}
 {\approx}$.
\end{lem}
\begin{proof}
నిరూపణ \olref[mod]{lem:val-in-termmodel} నిరూపణను పోలి ఉంటుంది.
\end{proof}

\begin{prob}
\olref[fol][com][ide]{lem:val-in-termmodel-factored} నిరూపణను పూర్తి
చేయండి.
\end{prob}

\begin{lem}
\ollabel{lem:truth}
అన్ని వాక్యాలు~$!A$కు $\Sat{\equivclass{M}{\approx}}{!A}$ అయినప్పుడు,
అప్పుడు మాత్రమే $!A \in \Gamma^*$.
\end{lem}

\begin{proof}
\olref[mod]{lem:truth} నిరూపణలోలాగే~$!A$పై ఆగమనం ద్వారా. $!A \ident
\eq[t][t']$ అయిన సందర్భం మాత్రమే అదనపు శ్రద్ధ కోరుతుంది.
\begin{align*}
\Sat{\equivclass{M}{\approx}}{\eq[t][t']} & \text{ iff } \equivrep{t}{\approx} = \equivrep{t'}{\approx}
\text{ (by definition of $\Struct{\equivclass{M}{\approx}}$)}\\
& \text{ iff } t \approx t' \text{ (by definition of $\equivrep{t}{\approx}$)}\\
& \text{ iff } \eq[t][t'] \in \Gamma^* \text{ (by definition of $\approx$).}
\end{align*}
\end{proof}

\begin{digress}
$\Struct{M(\Gamma^*)}$ ఎల్లప్పుడూ !!{enumerable}, అనంతమైనదే అయినా,
$\Struct{\equivclass{M}{\approx}}$ పరిమితంగా ఉండవచ్చు; ఎందుకంటే
$\equivrep{t}{\approx}$ వర్గాలు పరిమిత సంఖ్యలో మాత్రమే ఉన్నాయని
తేలవచ్చు. ఇది ఊహించదగినదే; $\Gamma$లో, వాటిని సత్యం చేసే ఏ
!!{structure} అయినా పరిమితంగా ఉండాలని కోరే !!{sentence}s ఉండవచ్చు.
ఉదాహరణకు, $\lforall[x][\lforall[y][x = y]]$ ఒక అవైరుధ్య
!!{sentence}; కానీ సరిగ్గా ఒక !!{element} మాత్రమే గల !!a{domain} ఉన్న
!!{structure}sలోనే అది సంతృప్తిపరచబడుతుంది.
\end{digress}

\end{document}
